%%=============================================================================
%% Voorwoord
%%=============================================================================

\chapter*{\IfLanguageName{dutch}{Woord vooraf}{Preface}}
\label{ch:voorwoord}

Het schrijven van deze bachelorproef vormt het sluitstuk van mijn opleiding Toegepaste Informatica aan HOGENT. Hoewel mijn opleiding als Full Stack Developer zich voornamelijk richt op het ontwikkelen van software, heb ik bewust gekozen voor een onderwerp rond artificiële intelligentie. De snelle opkomst van AI en de steeds grotere impact ervan op softwareontwikkeling maakten dit voor mij een bijzonder interessante en relevante piste om verder te verkennen.

\vspace{1em}

In deze bachelorproef werd onderzocht in welke mate Large Language Models (LLM’s) bruikbare en betrouwbare feedback kunnen genereren op de codekwaliteit van studentenprojecten binnen de opleidingsonderdelen Front-end Web Development en Web Services. Dit onderzoek vertrekt vanuit een concreet probleem, waarbij de evaluatie van projecten tijdsintensief is en onderhevig kan zijn aan inconsistenties. Door middel van een literatuurstudie, interviews met lectoren en de ontwikkeling van een proof-of-concept werd nagegaan hoe LLM’s hierin een ondersteunende rol kunnen spelen.

\vspace{1em}

Het uitvoeren van dit onderzoek was een leerrijk proces waarin ik niet alleen mijn technische vaardigheden verder heb ontwikkeld, maar ook kritisch heb leren nadenken over de rol van AI binnen het programmeeronderwijs. Het bood mij de kans om mijn interesse in softwareontwikkeling te combineren met de mogelijkheden van nieuwe technologieën.

\vspace{1em}

Graag wil ik een aantal personen bedanken die hebben bijgedragen aan deze bachelorproef. In de eerste plaats wil ik mijn promotor, Mevr. V. Depestele, bedanken voor de begeleiding, de waardevolle feedback en de ondersteuning doorheen het volledige traject. Daarnaast wil ik mijn co-promotor, Dhr. T. Aelbrecht, bedanken voor het meedenken en de technische inzichten die hebben bijgedragen aan dit onderzoek. Verder wil ik ook de lector(en) bedanken die hebben meegewerkt aan het onderzoek en hun inzichten hebben gedeeld over het evaluatieproces van studentenprojecten. Hun input was van grote waarde voor de uitwerking van deze bachelorproef.

\vspace{1em}

Tot slot wil ik mijn familie en vrienden bedanken voor hun blijvende steun en aanmoediging tijdens mijn opleiding. Met deze bachelorproef hoop ik een bijdrage te leveren aan het inzicht in de mogelijkheden en beperkingen van LLM’s binnen het programmeeronderwijs, en een basis te bieden voor verder onderzoek naar de integratie van AI in evaluatieprocessen.
